\section{Evaluation}

\textit{Note:} A \textbf{partial live key-cell round} on Amazon DynamoDB (\texttt{eu-west-1}) filled the primary W3/W4 $\times$ K1--K3 $\times$ capacity cells ($n{=}1$ block, 100{,}000 seeded orders). Artefacts: \texttt{results/batches.csv} (\texttt{data\_source=live}), \texttt{report/generated/cell\_summary.md}. W1/W2 cells, full $n{=}30$ replication, and ANOVA remain \textbf{unfilled}. Moto unit tests exercise plumbing only and are not used as cell fills.

\subsection{Designed Campaign Configuration}

The full design (from \texttt{config/experiment.yaml}) and the executed key-cell round:

\begin{itemize}
    \item \textbf{Full design:} 1M synthetic orders / 10K customers; K1--K3 $\times$ on-demand/provisioned $\times$ W1--W4; $n{=}30$
    \item \textbf{This round (live):} 100{,}000 orders; K1--K3 $\times$ both capacity modes $\times$ W3,W4; $n{=}1$ randomised block; Zipfian $s{=}1.070916$
    \item \textbf{Region / pricing:} \texttt{eu-west-1} list prices in \texttt{config/prices.yaml} (not a Cost Explorer cross-check)
    \item \textbf{Off-factorial:} K4 code paths excluded from CA2 evidence
\end{itemize}

\subsection{Live Key-Cell Results (W3 / W4)}

\subsubsection{Latency Distribution}

Table \ref{tab:sim-latency} reports measured mean and p95 latency (ms) from the live key-cell round.

\begin{table}[h]
\centering
\caption{Latency (ms) by Configuration and Workload --- live key-cell ($n{=}1$)}
\label{tab:sim-latency}
\begin{tabular}{|l|c|c|c|c|}
\hline
\textbf{Configuration} & \textbf{Mean (W3)} & \textbf{p95 (W3)} & \textbf{Mean (W4)} & \textbf{p95 (W4)} \\ \hline
K1-OnDemand & 4.344 & 5.838 & 4.553 & 5.848 \\ \hline
K1-Provisioned & 4.364 & 5.796 & 4.890 & 6.035 \\ \hline
K2-OnDemand & 4.482 & 5.883 & 4.803 & 5.876 \\ \hline
K2-Provisioned & 4.354 & 5.765 & 4.597 & 5.848 \\ \hline
K3-OnDemand & 5.010 & 6.116 & 5.246 & 6.265 \\ \hline
K3-Provisioned & 4.895 & 5.787 & 5.356 & 6.223 \\ \hline
\end{tabular}
\end{table}

\textit{Interpretation:} All twelve W3/W4 cells completed with zero throttles under the offered load. K3 means are slightly higher than K1/K2 (scatter-gather reads). These are single-replicate measurements; they do not replace the planned $n{=}30$ campaign.

\subsubsection{Throttling Behaviour}

\begin{table}[h]
\centering
\caption{Throttle Rate (\%) by Configuration and Workload --- live key-cell ($n{=}1$)}
\label{tab:sim-throttle}
\begin{tabular}{|l|c|c|}
\hline
\textbf{Configuration} & \textbf{Throttle Rate (W3)} & \textbf{Throttle Rate (W4)} \\ \hline
K1-OnDemand & 0.0 & 0.0 \\ \hline
K1-Provisioned & 0.0 & 0.0 \\ \hline
K2-OnDemand & 0.0 & 0.0 \\ \hline
K2-Provisioned & 0.0 & 0.0 \\ \hline
K3-OnDemand & 0.0 & 0.0 \\ \hline
K3-Provisioned & 0.0 & 0.0 \\ \hline
\end{tabular}
\end{table}

\subsubsection{Cost Analysis}

Table \ref{tab:sim-cost} applies \texttt{prices.yaml} list-price arithmetic to measured RCU/WCU (provisioned uses measured wall time $\times$ configured capacity floor). \textbf{No} Cost Explorer export is claimed.

\begin{table}[h]
\centering
\caption{Cost per 10,000 Operations (USD) --- live key-cell ($n{=}1$)}
\label{tab:sim-cost}
\begin{tabular}{|l|c|c|}
\hline
\textbf{Configuration} & \textbf{Cost (W3)} & \textbf{Cost (W4)} \\ \hline
K1-OnDemand & 0.00390 & 0.00389 \\ \hline
K1-Provisioned & 0.00233 & 0.00100 \\ \hline
K2-OnDemand & 0.00387 & 0.00390 \\ \hline
K2-Provisioned & 0.01050 & 0.00100 \\ \hline
K3-OnDemand & 0.00706 & 0.00706 \\ \hline
K3-Provisioned & 0.00713 & 0.00557 \\ \hline
\end{tabular}
\end{table}

\textit{Observed (this round only):} K3 on-demand cost/10k is higher than K1/K2 on-demand (higher RCU from scatter-gather). Provisioned W4 cost/10k is lower than on-demand for K1/K2 at this short wall duration; K2-provisioned W3 is an outlier at \$0.0105/10k because provisioned floor is charged over the batch wall even when request volume is modest.

\subsection{Statistical Analysis (insufficient replication)}

Two-way ANOVA requires multiple replicates per cell. This key-cell round has $n{=}1$; Table \ref{tab:sim-anova} therefore remains unfilled. Pipeline code in \texttt{analysis/analyse.py} is covered by fixture unit tests only.

\begin{table}[h]
\centering
\caption{ANOVA for Mean Latency (W3) --- awaiting $n{\ge}2$ live replicates}
\label{tab:sim-anova}
\begin{tabular}{|l|c|c|c|}
\hline
\textbf{Factor} & \textbf{F-statistic} & \textbf{p-value} & \textbf{Effect} \\ \hline
Key Design (K) & [TO BE FILLED] & [TO BE FILLED] & [TO BE FILLED] \\ \hline
Capacity Mode (C) & [TO BE FILLED] & [TO BE FILLED] & [TO BE FILLED] \\ \hline
Interaction (K $\times$ C) & [TO BE FILLED] & [TO BE FILLED] & [TO BE FILLED] \\ \hline
\end{tabular}
\end{table}

\subsection{What unit tests do and do not show}

Pytest against moto validates artefact plumbing (not CA2 empirical answers):

\begin{enumerate}
    \item Terraform modules encode six K1--K3 $\times$ mode tables; tags are project-only (no student name/ID).
    \item Workload generator Zipfian calibration and K1--K3 key schemas match design checks.
    \item Lambda handler records latency/throttle/capacity fields and counts throttles without inventing success.
    \item K3 scatter-gather path merges shard reads in unit tests.
    \item Cost helpers apply \texttt{prices.yaml} arithmetic --- \textbf{not} a live Cost Explorer cross-check.
    \item ANOVA/post-hoc code runs on fixtures; fixture significance is \textbf{not} a DynamoDB finding.
\end{enumerate}

Moto does \textbf{not} replicate network latency, adaptive/burst capacity, auto-scaling delay, or multi-tenant effects. Optional K4 code paths are excluded from this evidence narrative.

\subsection{Residual for CA2 evidence}

Live W3/W4 key-cells are filled ($n{=}1$, 100k seed). Remaining for COMPLETE: W1/W2 cells, full 1M seed / $n{=}30$ factorial, and ANOVA on multi-replicate live data. Practitioner guidance remains deferred until those residuals close.
